\chapter{Atlases}\label{ch:atlas}
\begin{goals}
\textbf{In this chapter} you turn one composed figure into a book: a page per bench, per section or per domain, in a single PDF, with the level in each caption, the key plan's line moving from page to page, and figure and page numbers running through — and you make it again next month by pressing \ui{Run}.
\end{goals}

A set of bench plans is forty figures that differ in one number. So is a set of cross-sections every 25 m, or the same plan filtered to each domain in turn. Nobody should lay those out forty times, and nobody should trust forty hand-made pages to be consistent with each other.

An \term{atlas} is one figure from the composer (\cref{ch:composer}) printed once per row of a table. The figure supplies everything that stays the same: the page, the frames, the legend, the title block. The table supplies what changes: a section position, a filter, a camera. The machinery is a routine (\cref{ch:routines}), so an atlas is a file you can read, rerun, and keep with the project.

\begin{figure}[htbp]\centering
\begin{adjustbox}{max width=\linewidth}\begin{tikzpicture}[font=\sffamily\footnotesize, every node/.style={text=ink}, >=Stealth,
  box/.style={draw=rule, fill=tint, rounded corners=1.5pt, line width=0.5pt, inner sep=5pt, align=center}]
\node[box] (fig) at (0,1.6) {the figure\\[1pt]\textcolor{muted}{\scriptsize page · frames · captions}};
\node[box] (tab) at (0,-0.4) {the table\\[1pt]\textcolor{muted}{\scriptsize level: 560, 580, 600 …}};
\node[box, fill=ambersoft, draw=amber] (step) at (3.8,0.6) {a \texttt{compo} step\\[1pt]\textcolor{muted}{\scriptsize once per row}};
\draw[->, rule, line width=0.8pt] (fig.east) -- (step.north west);
\draw[->, rule, line width=0.8pt] (tab.east) -- (step.south west);
% pages
\foreach \k/\l in {0/560, 1/580, 2/600} {
  \draw[rule, fill=white, line width=0.5pt] (6.3+\k*0.3,0.05-\k*0.14) rectangle ++(1.5,1.1);
}
\node[font=\sffamily\scriptsize, text=muted] at (7.65,0.3) {600 RL};
\draw[->, rule, line width=0.8pt] (step.east) -- (6.2,0.6);
\node[box] (book) at (10.6,0.6) {\texttt{atlas/plan.pdf}\\[1pt]\textcolor{muted}{\scriptsize one page per row}};
\draw[->, rule, line width=0.8pt] (8.5,0.6) -- (book.west);
\end{tikzpicture}\end{adjustbox}
\caption{An atlas. The figure is the constant, the table is the variable; each row makes a page, and the pages are bound into one PDF when the run ends.}\label{fig:atlas-flow}
\end{figure}

\section{A first atlas}\label{sec:atlas-first}

Build the figure first, as one good page. Put a 3D view on it showing the section you want to step — a plan slab for bench plans, a N–S or E–W section for cross-sections — add its legend, scale and north, and write a caption. Then press \ui{atlas…} in the composer's figure bar.

micro proposes a starter routine named after the figure, with \emph{atlas} appended, asks you to confirm the name, saves it in the project's \file{recipes/} folder, and opens it in the routine editor. Routines are files in a project, so the project must be saved first.

The starter makes reasonable guesses, all of which you can change:

\begin{itemize}
\item it steps along the axis of the view's own section: a plan steps in z, a N–S section in x, an E–W section in y;
\item it cuts the extent of everything loaded along that axis — the meshes and the topography included — into six equal parts, and takes the middle of each;
\item it gives each page a slab four-fifths of the spacing thick, so neighboring pages do not overlap;
\item it collects the pages into \file{atlas/}\emph{figure}\file{.pdf} at 200 dpi;
\item it hands each page its position as a value called \code{level}.
\end{itemize}

The status line tells you the one thing left to do by hand: write \code{\$\{level\}} in a caption, so each page says where it is. Then press \ui{Run}.

\begin{tryit}
\begin{enumerate}
\item Open the demo project and save it to a folder with \menu{File>Save project as…}.
\item Hide the topography, press \key{d} for a plan view and make it orthographic.
\item Turn on a plan section with a slab (\cref{ch:sections}).
\item Open \menu{File>Composer…}. Add \ui{Add 3D view}, \ui{Add legend}, \ui{Add scale bar}.
\item Add \ui{Add text} under the view and set it to\\
\code{**Figure \{fig\}.** Bench \$\{level\} m RL — \{count:deposit.csv:,\} blocks. Page \{page\} of \{pages\}.}
\item Rename the figure \code{bench plan} with \ui{rename}.
\item Press \ui{atlas…} and accept the name it offers. The routine editor opens on \emph{bench plan atlas}.
\item Press \ui{Preview} to see the six pages it will make, then \ui{Run}.
\item When the run finishes, open \file{atlas/bench\_plan.pdf} from the project folder: six pages, at 440, 526, 612, 698, 783 and 869 m RL, numbered 1 of 6 to 6 of 6. Each caption counts its own slab's blocks — 260,000 on the first page, 280,000 through the middle of the deposit, fewer on the last, where the topography has cut the model away.
\end{enumerate}
\end{tryit}

\shot{atlas-pages}{Pages 1, 3 and 6 of the demo's bench atlas side by side: the same layout, the slab stepping up through the deposit, and each caption giving its own level, its own block count and its page number. On page 6 the starter pit (gray) covers most of the bench; the topography is hidden.}

\section{Reading the routine}\label{sec:atlas-yaml}

The whole atlas is this:

\begin{lstlisting}[language=yaml]
name: "bench plan atlas"
tool: "routine"
params:
  steps:
    - each:
        level:
          - 440
          - 526
          - 612
          - 698
          - 783
          - 869
      steps:
        - tool: "compo"
          with:
            figure: "bench plan"
            view:
              section:
                axis: "z"
                position: "$level"
                thickness: 68.6
                clip: "slab"
            vars:
              level: "$level"
            book: "atlas/bench_plan.pdf"
            dpi: 200
\end{lstlisting}

That is what \ui{atlas…} writes on the demo; on your data the levels and the thickness are its guesses for your extent. Read it from the outside in. The \code{each} block holds a table — here one column, \code{level}, with six rows — and runs the steps inside it once per row. \code{\$level} is replaced by the row's value before the step runs. The one step inside is the \code{compo} tool: render a page of the figure \emph{bench plan}, with its view's section moved to this row's level, and add the page to the book.

To change the levels, edit the list — in the YAML, or as a grid in the routine editor. This is the first thing to edit: the guesses are spread through the bounding box of everything loaded, not over your benches. On the demo they are 86 m apart with slabs 68.6 m thick — six sections through the whole deposit rather than six benches. For true 10 m benches through the ore, list the bench midpoints and set the thickness to match:

\begin{lstlisting}[language=yaml]
    - each:
        level:
          - 575
          - 585
          - 595
          - 605
      steps:
        - tool: "compo"
          with:
            figure: "bench plan"
            view:
              section:
                axis: "z"
                position: "$level"
                thickness: 10
                clip: "slab"
            vars:
              level: "$level"
            book: "atlas/bench_plan.pdf"
            dpi: 200
\end{lstlisting}

Slabs of 10 m centered 10 m apart touch without overlapping, so each block lands on exactly one page, and the pages' block counts add up to the blocks between 570 and 610 m RL (\cref{sec:compo-statsscope}). Add a level per bench for the rest of the pit.

\section{What changes from page to page}\label{sec:atlas-view}

The \code{view} setting holds the changes applied to the figure's view frames for one page. There are three, and they combine.

\begin{table}[htbp]\small
\begin{tabularx}{\linewidth}{@{}lL@{}}\toprule
setting & meaning \\ \midrule
\code{section} & The section for this page. \code{axis} (\code{x}, \code{y} or \code{z}) or a \code{normal} of three numbers, or \code{dipDir} and \code{dip}; then \code{position}, \code{thickness}, and \code{clip} (\code{slab}, \code{front} or \code{behind}). The word \code{"off"} removes the section. \\
\code{camera} & The camera for this page: \code{look} (the bearing it looks toward), \code{elevation}, \code{distance}, \code{target} (three world coordinates), \code{ortho}, \code{zExag}. \\
\code{filter} & A filter for this page. One expression applies to every visible layer of the view; a map applies an expression per layer, by name. \\ \bottomrule
\end{tabularx}
\caption{The \code{view} setting of the \code{compo} tool.}\label{tab:atlas-view}
\end{table}

Whatever a page does not change, it takes from the figure: a page that sets only the section keeps the figure's camera, colors and filters.

\paragraph{A page per section.} For vertical sections, step \code{x} or \code{y}. If the camera in the figure looks square at the section, it still does on every page: the section moves, the direction of view does not.

\begin{lstlisting}[language=yaml]
- each:
    northing:
      - 7765200
      - 7765400
      - 7765600
      - 7765800
  steps:
    - tool: "compo"
      with:
        figure: "cross section"
        view:
          section:
            axis: "y"
            position: "$northing"
            thickness: 20
        vars:
          northing: "$northing"
        book: "atlas/sections.pdf"
\end{lstlisting}

\paragraph{A page per domain.} Nothing says an atlas must step a section. Step a filter instead, and the pages are the same plan, domain by domain:

\begin{lstlisting}[language=yaml]
- each:
    lito:
      - "HEMATITE"
      - "ITABIRITE"
      - "CANGA"
  steps:
    - tool: "compo"
      with:
        figure: "domain plan"
        view:
          filter:
            "deposit.csv": "LITO = \"$lito\""
        vars:
          lito: "$lito"
        book: "atlas/domains.pdf"
\end{lstlisting}

With a caption reading \code{\$\{lito\}: \{count:deposit.csv:,\} blocks, mean \{mean:FE:.1f\} \% Fe}, each page reports its own domain's count and mean, computed over what that page shows (inside its section, if the figure's view has one).

\paragraph{Two things at once.} An \code{each} inside an \code{each} is every combination — every domain on every bench. A single \code{each} with two columns of equal length pairs them row by row: section 1 with camera 1, section 2 with camera 2. The routine's \ui{Preview} lists every page before anything is rendered, so you can see which you wrote (\cref{ch:routines}).

\paragraph{A table from a file.} For a long or a maintained list — the mine's bench elevations, the section lines of a drilling grid — keep the table in the project as a CSV or a worksheet and name it under the \code{each}: \code{from: "benches.csv"}, as in the routine at the end of this chapter. Each column of the file becomes a \code{\$}name. The atlas then follows the file.

\section{What every page gets right}\label{sec:atlas-page}

Each page is the figure made again with that page's changes, and everything in the figure that depends on its view follows without being told:

\begin{description}
\item[Captions] \code{\$\{level\}} — or any value named under \code{vars} — is filled from the row. So is any expression over them: \code{\$\{level - 10\}} prints the toe of a 10 m bench. View fields such as \code{\{section\}}, \code{\{count\}} and \code{\{mean:FE\}} are read from this page's render — and count only what is inside this page's section, so each page reports its own bench, and the benches' block counts add up to the model's (\cref{sec:compo-statsscope}). Set the view's \ui{captions} row to \ui{whole model} for a caption that quotes the whole model on every page.
\item[Page and figure numbers] \code{\{page\}} and \code{\{pages\}} count through the book, and \code{\{fig\}} continues from one page to the next: Figure 1 on page 1, Figure 2 on page 2.
\item[Legends] are read again from the layers while the page's scene is on.
\item[Scale] A view with a locked scale stays at that scale on every page (\cref{sec:compo-drawing}). The scale bar, \code{\{scale\}} and the coordinate ticks follow each page's view.
\item[Key plans] The key plan keeps its own whole-model view; only the section line on it moves, to where this page's section is (\cref{sec:compo-keyplan}).
\item[Vector sections] are rebuilt for each page, so every page is line work (\cref{sec:compo-vector}).
\item[Charts and tables] on the figure are printed as they stand. They are not recomputed per page.
\end{description}

\begin{caution}
The last point is easy to miss. A grade–tonnage chart placed on an atlas figure shows the same curve on every page — the one the figure holds. Statistics that should change with the page belong in caption fields, which are computed per page, not in a chart frame.
\end{caution}

Your own 3D view is put back when the run ends, exactly as you left it, and the figure in the project is not changed: each page is a copy.

\section{One book, or a file per page}\label{sec:atlas-out}

A \code{compo} step sends its page to a \code{book}, to an \code{out} file, or to both (\cref{tab:compo-tool}). An atlas normally uses \code{book}: one PDF, its pages in the order of the rows, written once when the run's steps have finished. Use \code{out} when you want a file per page — a PNG of each bench for a slide deck — and put the row's value in its name, as the table shows.

A book carries its figure inside it, like any PDF from the composer, unless the figure's \ui{embed figure} is unticked (\cref{sec:compo-pdf}). Drop the atlas on micro a year later and the figure it was made from opens, ready to be run again.

Several books can come out of one routine: a step with \code{book: "atlas/plans.pdf"} in one \code{each} and another with \code{book: "atlas/sections.pdf"} in a second make two PDFs in one run.

Rendering is the slow part. Forty pages of a large model at 300 dpi take a while, since every view on every page is rendered to completion; 200 dpi is a sensible working resolution, and vector sections cost far less to print than rendered ones.

\section{Running it, and running it again}\label{sec:atlas-run}

An atlas runs like any routine (\cref{ch:routines}). \ui{Preview} lists each page and the file it goes to. The run opens a tracker showing each page as it is made, with \ui{Cancel} between pages. If a page fails — a layer the figure needs is not loaded, a filter names a column that is not there — the run stops at that step and says why, and you can retry it or skip it and carry on.

Afterwards the run's record lists each book with its page count and the SHA-256 of the file (\cref{ch:records}). If the routine has the no-overwrites guard, it stops before rendering anything when a book or page file it would write already exists.

Because an atlas is a routine, it sits naturally at the end of a longer one: check the model, recompute the columns, then print.

\begin{lstlisting}[language=yaml]
steps:
  - tool: "modelqa"
    with:
      layer: "deposit.csv"
    as: "mq"
  - check: "mq.errors = 0"
    say: "the model is clean before it is printed"
  - each:
      from: "benches.csv"
    steps:
      - tool: "compo"
        with:
          figure: "bench plan"
          view:
            section:
              axis: "z"
              position: "$level"
              thickness: 10
          vars:
            level: "$level"
            name: "$name"
          book: "atlas/bench_plans.pdf"
\end{lstlisting}

Next month, load the new model under the same layer name and run it again. Every page is rendered from the data as it then is; nothing in the figure or the routine needs to be touched. To print the same atlas for a different deposit, point the figure at the other data first with \ui{⇄ other data…} (\cref{sec:compo-retarget}), and give the routine that figure's name.

\begin{tip}
Design the figure on its hardest page. Step the section to the widest bench and the longest caption value before you press \ui{atlas…}: a layout that holds there holds everywhere, and a view with a locked scale and a fixed camera target keeps every page registered with the next, so a reader can flip through the book and watch the orebody change in place.
\end{tip}
